\subsection{Постановка задачи}

Проведённый обзор позволяет уточнить задачу работы. Apache Paimon --- lakehouse-система
с LSM-деревом, поддерживающая несколько файловых форматов; среди современных
колоночных форматов выделяется Vortex с нативным векторизованным декодированием,
специализированными кодеками (FastLanes, ALP, FSST) и проталкиванием предикатов, но
его поддержка в Apache Paimon отсутствует. Требуется определить, в каких сценариях
Vortex даёт выигрыш, реализовать его интеграцию и измерить характеристики.

Конкретизация задач.

\begin{enumerate}
    \item Разработка модуля интеграции. Реализовать модуль
        \texttt{paimon-vortex}: фабрики записи и чтения поверх нативной библиотеки
        Vortex (доступ через JNI), конвертер Paimon-предикатов в нативные выражения
        Vortex, формирование статистик записанных файлов для манифеста, поддержку
        проекции столбцов и многопоточного чтения.
    \item Функциональное тестирование. Проверить корректность артефакта
        локально в контейнерах: запись и чтение таблиц с первичным ключом и таблиц
        только для добавления, слияние версий, эволюция числа бакетов, согласованность
        с другими форматами на одинаковых данных.
    \item Развёртывание окружения нагрузочного тестирования. Средствами
        Ansible развернуть на кластере (Apache Flink, Apache Paimon, объектное
        хранилище) воспроизводимое окружение для запуска нагрузочных задач, сбора
        метрик и управления жизненным циклом кластера.
    \item Эксперименты. Провести две серии измерений: (а)~batch-аналитика
        на наборе TPC-DS масштаба \SI{100}{\giga\byte} (таблицы только для добавления,
        крупные файлы) --- сравнение Parquet и Vortex по времени выполнения запросов;
        (б)~near-real-time-сценарий --- таблица с первичным ключом, потоковая запись
        с захватом изменений, параллельные пакетные аналитические запросы во время
        записи и после неё --- сравнение Parquet, ORC и Vortex по пропускной
        способности записи и времени запросов разных типов. Результаты статистически
        обработать.
    \item Рекомендации. По результатам сформулировать рекомендации по выбору
        файлового формата хранения для различных классов нагрузок в Apache Paimon.
\end{enumerate}

Гипотезы, проверяемые в работе:

\begin{enumerate}
    \item на batch-аналитике с селективными предикатами и крупными файлами Vortex
        быстрее Parquet за счёт нативного проталкивания предикатов и
        специализированных кодеков;
    \item на потоковой near-real-time-нагрузке Vortex не уступает Parquet по
        пропускной способности записи;
    \item на запросах с агрегацией без селективного предиката Vortex проигрывает
        Parquet из-за накладных расходов границы JVM/native в отсутствие
        проталкивания агрегатов;
    \item поведение Vortex под конкурентным чтением во время потоковой записи
        отличается от Parquet/ORC ввиду различия в реализации доступа к объектному
        хранилищу.
\end{enumerate}

Критерии оценки --- время выполнения запросов (среднее по повторным запускам и
выборкам, с учётом разброса), пропускная способность записи (строк в секунду), число
успешно завершённых запросов под нагрузкой; для качественных выводов ---
устойчивость метрик (наличие или отсутствие катастрофических выбросов).
